\chapter{Axioms of Quantum Mechanics}
\label{aqm}

Formally a quantum system is defined by four axioms,
describing states,observables,dynamics, and probabilities.
\section{States}
    \label{states}

\begin{axiom}[States]
    \label{sts}
Physical state is a vector in complex separable Hilbert space $ \mathcal{H} $
\[ |\psi \rangle \in \mathcal{H} .\]
\end{axiom}
%description and discussion of states vectors rays etc

\section{Observables}
    \label{observables}

\begin{axiom}[Observables]
    \label{observs}
Observables are self-adjoint operators $ \hat A= \hat A^\dagger$ on $\mathcal H$. 
\end{axiom}
% discussion of observables SA HERM etc.

\section{Dynamics}
    \label{dynamics}
\begin{axiom}[Dynamics]
    \label{dyn}
    There is an observable $ \hat H$ (Hamiltonian),which generates time evolution
    (unitary group):
    \begin{equation}
       \label {evol}
 \hat U(t) = e^{-i \hat Ht/\hbar}.
    \end{equation}
\end{axiom}

    Evolution may be seen as action on observables (Heisenberg picture)
\begin{align}
   \hat A(t)&= \hat U(t) \hat A \hat U^{-1}(t), \label {heis1}\\
    \frac{d \hat A}{d t} &= \frac{i}{\hbar}[\hat H,\hat A].\label{heis2}
\end{align}
    Or as action on states (Schr\"odinger picture)
\begin{align}
    |\psi(t)\rangle&= \hat U^{-1}(t) |\psi \rangle, \label{schr1} \\
    \frac{d |\psi \rangle}{d t} &=-\frac{i}{\hbar} \hat H |\psi \rangle. \label{schr2}
\end{align}
%here about dynamics

\section{Probabilities}
\label{probabilities}
\begin{axiom}[Born rule]
    \label{born-rule}
    \[  B_\psi=\frac{\langle{\psi}|{\phi}\rangle}{\|{\psi}\| \|{\phi}\|}  \]
     is probability amplitude of transition between states. Corresponding 
     probability is $ P_\psi=|B_\psi|^2 $.
    Expectation value of an observable in a state 
    \begin{equation}
       \label {exp1}
    \langle \hat A \rangle_\psi = \frac{\langle \psi| \hat A |\psi \rangle} {\langle \psi|\psi \rangle}
    \end{equation}
    comprises all available information.
\end{axiom}
%here about prob,measurement, projection etc

